\section{Introduction}
\label{sec:intro}

In August 2024 a system called The AI Scientist produced machine-learning papers end to end
and reviewed them with an automated reviewer \citep{lu2024aiscientist}. Its successor
submitted three manuscripts to a peer-reviewed workshop, one of which cleared the human
acceptance threshold \citep{yamada2025aiscientistv2}. By mid-2026 the same idea appears in
governed hospital deployments running genome-wide association studies over 286{,}422
patient records \citep{huang2026nais}, in condensed-matter pipelines that go from an
11{,}083-paper corpus to a physics manuscript across 47 fresh-context sessions
\citep{huang2026grounded}, and in agentic ``scientific operating systems'' that compile
research intent into staged execution plans \citep{zheng2026scion}.

The literature describing these systems is now large enough that nobody reads all of it.
That is the first thing this paper supplies: a systematic map of \nCorpus{} papers on
autonomous AI research systems from January 2024 to July 2026, coded on ten dimensions and
released with its data, code, and decision log.

But mapping capability is the easy part, and it is not the interesting question. A parallel
thread in this same literature has been asking a harder one. \citet{takahara2026hep} observe
that in current research agents the hypothesis--test--belief cycle is ``buried in
unstructured logs, and no mechanism lets the agent or the human researcher audit that
process.'' \citet{li2026calibration} argues that the central question is no longer whether
systems can generate hypotheses or manuscripts but ``whether their scientific claims are
calibrated to the evidence that supports them,'' proposing that evidence should
\emph{license} claims. \citet{bowkis2026automated} show that automating research produces
systematic, hard-to-detect errors even when no agent is scheming, because optimisation
pressure concentrates mistakes precisely where human reviewers are least likely to catch
them. And \citet{gyevnar2026ddos} demonstrate the consequence at scale: a remote adversary
poisoning an open dataset compromised autonomous research agents in 49.6\% of runs while
being detected in 6\%.

So we asked the question that this thread implies but nobody had measured across the field:
\emph{when these systems make claims, what evidence do they offer that the claims are true?}
Not how good the systems are --- what they let you check.

The answer is lopsided. The field grew roughly \growthFactor{}-fold in nine quarters
(\firstQn{} papers in \firstQ{} to \peakQn{} in \peakQ{}), and by conventional standards it
evaluates itself diligently: \pctBenchMetric{} of the corpus reports benchmark metrics. But
only \pctHeldOut{} of papers re-test a selected result outside the loop that produced it,
\pctAuditNone{} provide no auditability mechanism of any kind --- no inspectable traces, no
claim-to-evidence provenance, no reproduction artifacts, no calibration --- and
\pctHumanUnspec{} never state what the humans did. Among the \nDiscovery{} papers claiming
the AI produced new scientific knowledge, \pctDiscNoAudit{} provide no auditability
mechanism. Benchmark scores are abundant; the ability to check a specific claim is scarce.

\subsection*{The instrument is also the object}

This review was conducted by a frontier language model --- Claude Fable 5 --- working under
a human researcher's direction. That is not a stunt, and it is not incidental to the
findings. A review of autonomous AI research, performed by an autonomous AI researcher, can
either quietly enact the problem it describes or document it. We chose to document it.

Every step was logged as it happened: a contemporaneous timeline, an error log of the AI's
own failures, and a record of every point where the human co-author intervened. Nothing was
reconstructed afterwards, and nothing was removed for looking bad. That log turns out to be
the most useful thing we produced, because it makes measurable what the corpus we studied
mostly leaves invisible:

\begin{itemize}\itemsep2pt
\item Two frontier models screening the same \nScreened{} papers under the same written
protocol agreed on only \rawAgree{} of decisions ($\kappa = \kappaThree{}$), and the
disagreements were not noise. They clustered on seven boundary classes --- assistive tools,
automated literature review, self-driving-lab infrastructure, human-AI-usage studies,
data-science agents, general-purpose deep-research agents, and research-ethics frameworks ---
where the field has not settled what counts as an AI doing science. The protocol, written by
the AI and approved by the human, looked rigorous and left a fifth of its own decisions
undetermined.

\item A workflow reported success in 106 milliseconds having done nothing at all. Structured
extraction silently dropped items at a measured rate of 0.4--0.5\%. A pipeline defect fed
empty inputs to a reasoning step, which duly reasoned about nothing --- and it hit precisely
the two records that human expertise had rescued from the search queries.

\item The human co-author overruled the AI's most consequential scope judgment, and was
right to; that single reversal changed the corpus by 183 papers.
\end{itemize}

None of these failures would be visible in a finished paper. Each was caught by a cheap
mechanical check --- expected-versus-actual reconciliation, refuse-on-empty-input, running
two raters against each other --- that costs almost nothing and that the great majority of
the corpus we mapped does not report performing. That is the argument this paper ends on:
the auditability deficit we measure in the literature and the failure modes we logged in
ourselves are the same problem, and it is tractable.

\subsection*{Contributions}

\begin{enumerate}\itemsep2pt
\item A systematic map of \nCorpus{} papers (2024--2026) on autonomous AI research systems,
coded on ten dimensions, with full data, code, protocol, and decision log released.
\item A measurement of the evidentiary practices of that literature --- what is evaluated,
what is auditable, what is released --- including a mechanical full-text verification of
artifact release on \nArtChecked{} papers, which we show corrects a substantial bias in
abstract-based coding (\pctRepoAbstract{} of papers state code release in the abstract;
\pctRepo{} actually link a repository).
\item A reflexive case study of AI-conducted systematic review, with measured inter-model
reliability, a documented protocol failure at the field's conceptual boundaries, and an
error taxonomy with observed rates.
\item A concrete claim: the checks that caught our failures are cheap, and their near-absence
from the literature is a solvable reporting problem rather than a hard research problem.
\end{enumerate}
